\documentclass[12pt,a4paper]{report}

\usepackage{iftex}
\ifPDFTeX
\usepackage[utf8]{inputenc}
\usepackage[T1]{fontenc}
\newcommand{\malayalamexample}{\textit{myre in Malayalam script}}
\else
\usepackage{fontspec}
\setmainfont{Latin Modern Roman}
\newfontfamily\malayalamfont{Noto Sans Malayalam}
\newcommand{\malayalamexample}{{\malayalamfont മൈര്}}
\fi
\usepackage{amsmath}
\usepackage{array}
\usepackage{booktabs}
\usepackage{geometry}
\usepackage{graphicx}
\usepackage{longtable}
\usepackage{setspace}
\usepackage{xcolor}
\usepackage{listings}
\usepackage{fancyhdr}
\usepackage{hyperref}
\hypersetup{hidelinks}
\usepackage{lipsum}

\geometry{a4paper, margin=1in}
\onehalfspacing
\setlength{\parindent}{0.5in}
\setlength{\parskip}{0.5em}
\setlength{\headheight}{30pt}
\setlength{\headsep}{18pt}
\sloppy

\newcommand{\projecttitle}{Safe Chat: Real-Time Manglish \& Malayalam Hate Speech Detection System}
\newcommand{\shortprojecttitle}{Safe Chat: Manglish \& Malayalam}
\newcommand{\authorname}{Bastin George}
\newcommand{\organization}{ICFOSS (International Centre for Free and Open Source Software)}
\newcommand{\reportmonth}{\today}
\newcommand{\headerfont}{\normalfont\scriptsize}
\newcommand{\footerfont}{\normalfont\footnotesize}

\renewcommand{\chaptermark}[1]{\markboth{#1}{}}
\renewcommand{\sectionmark}[1]{}

\newcommand{\leftpageheader}{%
    \parbox[b]{0.48\headwidth}{\raggedright\headerfont\bfseries\shortprojecttitle}%
}

\newcommand{\rightpageheader}{%
    \parbox[b]{0.48\headwidth}{\raggedleft\headerfont\nouppercase{\leftmark}}%
}

\newcommand{\alignedpagefooter}{%
    \begin{tabular*}{\headwidth}{@{\extracolsep{\fill}}l r@{}}
        {\footerfont\organization} & {\footerfont Page~\thepage}
    \end{tabular*}%
}

\fancypagestyle{alignedreport}{%
    \fancyhf{}%
    \fancyhead[L]{\leftpageheader}%
    \fancyhead[R]{\rightpageheader}%
    \fancyfoot[C]{\alignedpagefooter}%
    \renewcommand{\headrulewidth}{0.4pt}%
    \renewcommand{\footrulewidth}{0.4pt}%
}

\fancypagestyle{plain}{%
    \fancyhf{}%
    \fancyhead[L]{\leftpageheader}%
    \fancyhead[R]{\rightpageheader}%
    \fancyfoot[C]{\alignedpagefooter}%
    \renewcommand{\headrulewidth}{0.4pt}%
    \renewcommand{\footrulewidth}{0.4pt}%
}

\pagestyle{alignedreport}

\lstdefinestyle{codeStyle}{
    basicstyle=\ttfamily\scriptsize,
    breaklines=true,
    frame=single,
    numbers=left,
    numberstyle=\tiny\color{gray},
    keywordstyle=\color{blue},
    commentstyle=\color{purple},
    stringstyle=\color{teal},
    showstringspaces=false,
    tabsize=4
}
\lstset{style=codeStyle}

\newcommand{\missingfigure}[1]{%
    \fbox{%
        \parbox[c][2in][c]{0.8\textwidth}{%
            \centering\textit{Image file \texttt{\detokenize{#1}} is not available in this workspace.}%
        }%
    }%
}

\newcommand{\safeincludegraphics}[2][]{%
    \IfFileExists{#2}{\includegraphics[#1]{#2}}{\missingfigure{#2}}%
}

\newcommand{\missinglisting}[1]{%
    \par\noindent
    \fbox{%
        \parbox{0.9\textwidth}{%
            \textit{Source file \texttt{\detokenize{#1}} is not available in this workspace.}%
        }%
    }%
    \par
}

\newcommand{\safeinputlisting}[2][]{%
    \IfFileExists{#2}{\lstinputlisting[#1]{#2}}{\missinglisting{#2}}%
}

\renewcommand{\contentsname}{CONTENTS}
\renewcommand{\listfigurename}{LIST OF FIGURES}
\renewcommand{\listtablename}{LIST OF TABLES}
\renewcommand{\bibname}{REFERENCES}

\begin{document}

\pagenumbering{gobble}

\begin{titlepage}
\begin{center}
\vspace*{1.5cm}
{\Large \textbf{SAFE CHAT: REAL-TIME MANGLISH \& NATIVE MALAYALAM HATE SPEECH DETECTION SYSTEM}}\\[1.5cm]
{\large \textbf{TECHNICAL PROJECT REPORT}}\\[1.5cm]
submitted by\\[0.8cm]
{\large \textbf{\authorname{} }}\\[1.5cm]
prepared for\\[0.8cm]
{\Large \textbf{\organization{}}}\\[0.5cm]
\vfill
{\large \reportmonth{}}
\end{center}
\end{titlepage}

\clearpage
\pagenumbering{roman}
\setcounter{page}{1}

\chapter*{ABSTRACT}
\markboth{ABSTRACT}{}
\addcontentsline{toc}{chapter}{ABSTRACT}
The rapid and unprecedented expansion of digital communication channels in the state of Kerala has fundamentally transformed how individuals interact online. This digital revolution has led to the absolute dominance of Manglish—the practice of writing the Malayalam language using the English alphabet—as well as the continued heavy usage of Native Malayalam script for informal, text-based interaction. While this linguistic evolution has bridged digital divides, it has also created a catastrophic vulnerability in online safety. The Safe Chat System is an advanced, real-time messaging application designed to solve the critical issue of automated hate speech, cyberbullying, and profanity detection in highly diverse, code-mixed linguistic environments, a domain where standard, commercially available Natural Language Processing (NLP) models inherently fail.

This project introduces a highly sophisticated Artificial Intelligence architecture centered around a Hugging Face MuRIL (Multilingual Representations for Indian Languages) model, aggressively fine-tuned on a colossal, custom-engineered dataset of nearly 40,000 unique records. To combat the severe limitations of standard AI—specifically dataset bias, aggressive trolling techniques (such as character dragging and deliberate misspellings), and the ability of Native script to bypass English-biased tokenizers—the system incorporates an impenetrable, multi-layered mathematical pre-processing firewall. 

This firewall features a dedicated Transliteration Layer utilizing the \texttt{indic-transliteration} library to automatically normalize Native Malayalam inputs into OPTITRANS Manglish prior to classification, effectively closing the native script vulnerability. It further utilizes complex Regular Expressions (Regex) for surgical semantic disambiguation—such as isolating food contexts (e.g., "pooryum curryum") to prevent devastating false positive classifications on identical-sounding words—and a TF-IDF Cosine Similarity engine that catches infinite spelling variations of slurs without relying on rigid keyword blacklists.

Furthermore, the Safe Chat System pioneers a completely live, asynchronous chat room experience driven by an aggressive frontend polling mechanism, mimicking the responsiveness of modern commercial messaging applications like WhatsApp and Telegram. It introduces a highly efficient Speech-to-Text (STT) dictation pipeline utilizing the browser's native Web Speech API, allowing users to transcribe their voice directly into the messaging portal in real-time, bypassing the heavy processing costs associated with server-side models like OpenAI's Whisper.

Finally, with a continuous active learning loop and an Instant Whitelist Cache, the platform allows administrative UI feedback to dynamically update the core dataset and correct the AI's contextual biases on the fly. The result is a robust, self-healing, and highly accurate moderation ecosystem capable of securing regional digital communication channels.

\clearpage
\tableofcontents
\listoffigures

\clearpage
\pagenumbering{arabic}
\setcounter{page}{1}

\chapter{INTRODUCTION}

\section{Background of the Study}
In the modern era, digital communication has become an indispensable facet of academic, personal, and professional life. The proliferation of affordable smartphones and ubiquitous internet connectivity has brought millions of users online, leading to an explosion of user-generated content on social media platforms, forums, and instant messaging applications. In regions characterized by rich linguistic heritage, such as Kerala, India, this digital shift has given rise to unique linguistic phenomena. Users heavily rely on "Manglish"—a code-mixed language that combines Malayalam vocabulary and grammar with the English (Latin) script. Furthermore, the use of Native Malayalam script remains highly prevalent, often used interchangeably with Manglish in the same conversation.

Because Manglish is entirely informal, it possesses no standardized spelling rules, no formal grammar constraints, and relies heavily on phonetic typing. A single word can be spelled in dozens of different ways depending on the user's dialect, typing speed, and personal preference. Consequently, traditional Artificial Intelligence (AI) and Natural Language Processing (NLP) filters, which are predominantly trained on structured, monolingual English datasets, are fundamentally incapable of monitoring these conversations effectively. 

\section{The Threat of Cyberbullying and Hate Speech}
The anonymity and reach of the internet have exacerbated issues of unsafe online communication, cyberbullying, targeted harassment, and hate speech. The lack of effective moderation tools for regional and code-mixed languages has turned many regional online spaces into toxic environments. When users can harass others using regional slurs spelled phonetically in English, standard safety algorithms fail to trigger. This leaves vulnerable populations exposed to abuse and creates a pressing demand for localized, intelligent moderation systems.

\sectionsection{Problem Statement}
Standard messaging platforms either lack automated moderation entirely or rely on rudimentary, hardcoded keyword-blocking systems. These legacy systems are easily bypassed by malicious users (trolls). A user can simply misspell a slur (e.g., typing a slur as "myyyyyyrr" or "m@yre" instead of its base form) to evade detection. Furthermore, a troll can type the slur in Native Malayalam script, completely bypassing western AI models that only understand ASCII characters. 

Additionally, machine learning models suffer from severe dataset bias. They often misclassify innocent words simply because those words share substrings with vulgar terms, or because they frequently appear in abusive contexts within the training data. For example, the Malayalam word for a specific food item ("poori") shares its exact spelling with a severe derogatory slur. A standard AI will flag discussions about food as hate speech, ruining the user experience. There is a critical need for a live chat ecosystem that utilizes advanced mathematical pre-processing to normalize text, a robust transformer-based AI to understand deep context, and a live feedback loop to correct biases on the fly.

\section{Objectives}
The primary objective of the Safe Chat System is to research, design, and deploy an intelligent, adaptive messaging platform capable of securing code-mixed Malayalam communication. The specific objectives are defined as follows:
\begin{itemize}
    \item \textbf{Data Engineering:} To engineer a massive, robust code-mixed dataset by merging, mapping, and mathematically deduplicating data from multiple open-source repositories and synthetic generation scripts, resulting in a balanced training foundation.
    \item \textbf{Pre-processing Firewall:} To build an advanced pre-processing firewall featuring Transliteration, Regex contextual disambiguation, and TF-IDF Cosine Similarity to normalize exaggerated text, catch spelling variations, and prevent false positives.
    \item \textbf{AI Fine-Tuning:} To fine-tune the state-of-the-art Hugging Face MuRIL architecture, teaching it to accurately classify code-mixed text into distinct, nuanced safety buckets rather than a simple binary output.
    \item \textbf{Live Asynchronous Architecture:} To implement a Live Polling Architecture on the frontend, establishing a realistic, asynchronous chat experience that mimics commercial platforms without overwhelming server resources.
    \item \textbf{Voice Integration:} To integrate a Speech-to-Text (STT) dictation pipeline utilizing the lightweight Web Speech API, enabling seamless Malayalam voice typing.
    \item \textbf{Active Learning:} To establish a continuous active learning pipeline that captures user feedback from the UI and automatically re-trains the model upon server shutdown, ensuring the AI evolves alongside internet slang.
\end{itemize}

\chapter{SYSTEM REQUIREMENTS AND SPECIFICATIONS}
\section{Hardware Requirements}
The development and deployment of the Safe Chat AI model requires substantial computational resources, primarily due to the heavy demands of Transformer-based NLP architectures.
\begin{itemize}
    \item \textbf{Processor (CPU):} Minimum Intel Core i5 or AMD Ryzen 5. An 8-core processor is recommended for parallel data processing.
    \item \textbf{Memory (RAM):} Minimum 16 GB RAM. 32 GB is highly recommended for loading large datasets into memory during the data engineering phase.
    \item \textbf{Graphics Processing Unit (GPU):} A dedicated NVIDIA GPU (e.g., RTX 3060, RTX 4070, or cloud equivalents like Tesla T4/V100) with at least 8 GB of VRAM is absolutely necessary for fine-tuning the MuRIL model within a practical timeframe. 
    \item \textbf{Storage:} 50 GB of fast NVMe SSD storage for datasets, virtual environments, and storing large model weight files (\texttt{model.bin} and \texttt{.safetensors}).
\end{itemize}

\section{Software Requirements}
The system is built heavily on Python, leveraging the most modern data science and web frameworks available.
\begin{itemize}
    \item \textbf{Operating System:} Windows 10/11, macOS, or any modern Linux distribution (Ubuntu 20.04+ recommended for server deployment).
    \item \textbf{Language:} Python 3.9 or higher.
    \item \textbf{Deep Learning Frameworks:} \texttt{PyTorch} and Hugging Face \texttt{Transformers} for defining and training the neural network.
    \item \textbf{Data Engineering Libraries:} \texttt{Pandas} and \texttt{NumPy} for massive dataset manipulation; \texttt{scikit-learn} for machine learning metrics and TF-IDF implementation.
    \item \textbf{Natural Language Processing Libraries:} \texttt{indic-transliteration} for converting Native script to Manglish.
    \item \textbf{Backend Server:} \texttt{Flask} for the RESTful API and rendering HTML templates.
    \item \textbf{Frontend Stack:} Vanilla JavaScript, HTML5, CSS3. Specifically, utilizing the Web Speech API natively built into modern browsers like Google Chrome and Mozilla Firefox.
\end{itemize}

\chapter{SYSTEM ARCHITECTURE AND DATA ENGINEERING}

\section{Data Engineering and Dataset Construction}
The foundation of any robust AI system is its training data. For the Safe Chat System, the dataset required meticulous engineering due to the lack of pre-existing, high-quality Manglish hate speech repositories. 

The dataset was constructed by parsing and merging eight separate massive datasets. A primary source was the \textit{DravidianCodeMix-2020} dataset, which provided a baseline of offensive and sentiment-annotated data. However, sentiment data (e.g., categorizing text as Negative or Positive) does not directly map to offensive categories. A negative sentiment (e.g., "I am very sad today") is not hate speech. Therefore, a custom mathematical label mapping function was built via Python scripts to translate non-abusive negative sentiment into safe buckets, ensuring the model would not confuse negativity with toxicity.

To counter severe class imbalance—where safe messages vastly outnumbered hate speech—synthetic data generation scripts were deployed. These scripts generated thousands of variations of known slurs and offensive phrases. Following a synthetic Cartesian expansion, the final master dataset (\texttt{mal\_full\_offensive\_train.csv}) was compressed and mathematically deduplicated, resulting in exactly 39,716 completely unique, highly optimized records. This dataset serves as the definitive textbook for the AI model.

\section{The Multi-Layered Preprocessing Firewall}
Before user input is ever passed to the neural network for classification, it is intercepted by a highly strict, multi-layered pre-processing firewall. This firewall is designed to catch trolls, normalize spelling variations, and correct contextual dataset bias, ensuring the AI only processes clean, disambiguated text.

\subsection{The Transliteration Layer (Native Malayalam Support)}
One of the most significant architectural vulnerabilities discovered during development was the system's susceptibility to Native Malayalam script. Trolls could effortlessly bypass the hate speech filters by typing slurs in their native script (e.g., typing "\malayalamexample{}" instead of "myre"). Because the core training dataset was heavily biased towards English-script Manglish, the model failed to recognize the Native script equivalents, scoring severe hate speech as 99\% safe.

To solve this critical vulnerability without the need to source, annotate, and train a completely new Native dataset, a \textbf{Transliteration Layer} was engineered inside the Backend firewall. Leveraging the \texttt{indic-transliteration} Python package, all incoming messages are passed through the \texttt{sanscript} module. Any Native Malayalam characters detected are instantly and mathematically converted into OPTITRANS (ASCII-based Manglish) before the AI model ever sees them. This normalization ensures that a single AI model can effectively moderate both Manglish and Native Malayalam simultaneously, providing a unified defense mechanism.

\subsection{Regex Contextual Disambiguation}
Following transliteration, the text enters a sophisticated regular expression (Regex) engine. This engine performs two critical functions:
\begin{enumerate}
    \item \textbf{Character Reduction:} It strips exaggerated typing and character dragging. A troll attempting to bypass filters by typing "myyyyyyreeeeeee" will have their input mathematically reduced to its base phonetic structure.
    \item \textbf{Semantic Isolation (The Food Context Fix):} The Regex engine surgically separates identical-looking words based on semantic context, solving severe dataset bias. A major achievement of this layer is Contextual Protection. For example, the Malayalam word "poori" translates to a puffed bread (a food item), but shares its exact spelling with a severe derogatory slur. Standard AI models inevitably flag discussions about food as hate speech. The Regex firewall is programmed to detect surrounding food-related vocabulary (e.g., "curry", "bhaji", "masala") and instantly neutralize the phrase by replacing the entire substring with a safe token (e.g., "food") before classification. This entirely eliminates false positive predictions, protecting normal user conversations.
\end{enumerate}

\subsection{Mathematical Foundations: TF-IDF Cosine Similarity Engine}
The final layer of the firewall is the Cosine Similarity Engine. Using \textit{scikit-learn}, character n-grams (groupings of 2 to 3 characters) are extracted from the user's text and mapped to mathematical vectors. The engine relies on two foundational concepts:

\textbf{1. Term Frequency-Inverse Document Frequency (TF-IDF):}
This numerical statistic reflects how important a word or n-gram is to a document in a collection or corpus.
The Term Frequency (TF) measures how frequently a term $t$ appears in a document $d$:
\begin{equation}
TF(t, d) = \frac{\text{Number of times } t \text{ appears in } d}{\text{Total number of terms in } d}
\end{equation}

The Inverse Document Frequency (IDF) measures how important a term is across the entire corpus $D$:
\begin{equation}
IDF(t, D) = \log \left( \frac{\text{Total number of documents } N}{\text{Number of documents containing term } t} \right)
\end{equation}

The TF-IDF weight is the product of these two values:
\begin{equation}
\text{TF-IDF}(t, d, D) = TF(t, d) \times IDF(t, D)
\end{equation}

\textbf{2. Cosine Similarity:}
Once the text is converted into vectors, we measure the cosine of the angle between them to determine similarity, irrespective of their magnitude. For vectors $A$ and $B$:
\begin{equation}
\text{Cosine Similarity}(A, B) = \cos(\theta) = \frac{A \cdot B}{\|A\| \|B\|} = \frac{\sum_{i=1}^{n} A_i B_i}{\sqrt{\sum_{i=1}^{n} (A_i)^2} \sqrt{\sum_{i=1}^{n} (B_i)^2}}
\end{equation}

Operating on a strict 0.85 threshold, the engine evaluates phonetic and structural similarity. If a user types a highly obfuscated variation of a slur that evades Regex, the Cosine Similarity engine will mathematically determine its proximity to the base slur vector and automatically snap it to the known offensive token. This approach catches infinite spelling variations.

\chapter{ARTIFICIAL INTELLIGENCE AND FINE-TUNING}

\section{The MuRIL Architecture and Transformer Mechanics}
The core intelligence of the Safe Chat System is powered by \textbf{MuRIL} (Multilingual Representations for Indian Languages), developed by Google. MuRIL is a BERT-based transformer model pre-trained on 17 Indian languages and their transliterated counterparts.

\subsection{Self-Attention Mechanism}
At the heart of the MuRIL model is the Self-Attention mechanism, which allows the model to weigh the importance of different words in a sentence relative to each other, capturing deep semantic context. 
For a given sequence of word embeddings, the model creates three matrices: Query ($Q$), Key ($K$), and Value ($V$). The attention score is calculated as:
\begin{equation}
\text{Attention}(Q, K, V) = \text{softmax}\left(\frac{QK^T}{\sqrt{d_k}}\right)V
\end{equation}
where $d_k$ is the dimension of the key vectors. This equation mathematically explains how MuRIL understands that the word "myre" acts differently depending on the surrounding words in the sentence.

\section{The Fine-Tuning Pipeline}
The fine-tuning process was orchestrated using the Hugging Face \texttt{transformers} library. The custom dataset was tokenized and fed into the MuRIL architecture using a \texttt{SequenceClassification} head. 

The model was trained to classify text into five nuanced safety buckets, moving beyond a simple binary "Safe/Unsafe" output. The labels include:
\begin{itemize}
    \item \textbf{Not\_offensive:} Completely safe, normal chat.
    \item \textbf{Profanity:} General vulgarity and swear words without a specific target.
    \item \textbf{Off\_target\_ind:} Hate speech or abuse targeted at a specific individual.
    \item \textbf{Off\_target\_group:} Hate speech targeted at a specific group (e.g., religious, political, or gender-based).
    \item \textbf{Not\_in\_intended\_language:} Text that the model cannot process or falls outside the domain.
\end{itemize}

During training, the optimization process utilized the Cross-Entropy Loss function to minimize classification error:
\begin{equation}
\text{Loss} = -\sum_{c=1}^{M} y_{o,c} \log(p_{o,c})
\end{equation}
where $M$ is the number of classes (5), $y_{o,c}$ is a binary indicator of whether class $c$ is the correct classification for observation $o$, and $p_{o,c}$ is the predicted probability that observation $o$ belongs to class $c$. 

The model achieved remarkable accuracy, mapping complex morphological structures of Manglish to their offensive categories. The final weights were saved into the \texttt{finetuned\_model/} directory.

\chapter{RECENT IMPLEMENTATIONS AND CAPABILITIES}

\section{Telegram Bot Integration and Native Censor Bars}
To extend the reach of the Safe Chat system beyond the web UI, a robust Telegram Bot (\texttt{telegram\_bot.py}) was engineered. Utilizing the \texttt{python-telegram-bot} library, the bot actively monitors group chats, intercepting both text messages and voice notes in real-time. Instead of using native Telegram "Spoilers" (which users can simply tap to reveal), the bot pioneers a "Permanent Censor Bar" mechanism. When toxic text is detected, the bot instantly deletes the malicious message and programmatically replaces the exact length of the original sentence with block characters (e.g., \texttt{████████}). This completely sanitizes the chat environment while retaining the conversation's flow.

\section{Continuous Reinforcement Learning Loop}
To prevent the model from stagnating as internet slang evolves, a continuous active learning loop was established. The web UI features a "Report as Hate Speech" button, and an administrative Labeling Tool allows humans to correct the AI's predictions via a \texttt{/api/classify} endpoint. These corrections are actively logged and injected back into the master dataset (\texttt{mal\_full\_offensive\_train.csv}). 

To automate the retraining phase, an \texttt{update\_brain.py} script was developed and tethered to the Flask server's shutdown hook. When the server is gracefully terminated, the system automatically detects new user feedback, splits a fresh 80/20 train/test holdout, and initiates a Full Parameter Fine-Tuning sequence, allowing the AI to effectively "sleep and study" before waking up smarter.

\section{Whisper Speech-to-Text (STT) Voice Processing}
Voice notes are a primary medium of communication in modern chat apps. To moderate audio, the system integrates the OpenAI \texttt{Whisper} library directly into the Flask backend. When a user sends a voice note via the Web UI or Telegram, the audio is converted via \texttt{FFmpeg}, mathematically transcribed by Whisper, and the resulting Malayalam transcript is fed seamlessly into the AI text pipeline for toxicity classification.

\section{Deep Slang Preprocessing and Neutralization}
To combat the AI's tendency to falsely flag casual conversation, the \texttt{preprocessing.py} firewall maps heavily used, harmless Manglish slang (such as "da", "poda", "aliyan", "ninte") into the neutral mathematical token "friend". This dramatically reduces false positives and ensures the AI focuses strictly on malicious intent rather than casual linguistic tone.

\chapter{RESULTS AND EVALUATION}

\section{Evaluation Metrics}
The performance of the fine-tuned MuRIL model was rigorously tested using an array of evaluation scripts (\texttt{evaluate\_strict.py}, \texttt{evaluate\_metrics.py}). The model's predictions were compared against a holdout testing dataset to calculate Accuracy, Precision, Recall, and the F1-Score.

\begin{itemize}
    \item \textbf{Accuracy:} The percentage of total predictions that were correct.
    \item \textbf{Precision:} The ability of the classifier not to label a negative sample as positive (minimizing false positives).
    \item \textbf{Recall:} The ability of the classifier to find all the positive samples (minimizing false negatives).
    \item \textbf{F1-Score:} The harmonic mean of Precision and Recall, providing a single metric for model performance, especially crucial for imbalanced datasets.
\end{itemize}

\section{Visualizations}
The evaluation pipeline automatically generated several visual diagnostic tools to aid in understanding the model's behavior:
\begin{enumerate}
    \item \textbf{Confusion Matrix (\texttt{confusion\_matrix\_unseen.png}):} This heatmap visualization provided a clear breakdown of true positives versus false positives across all five classification buckets. It highlighted areas where the model struggled, such as distinguishing between general profanity and targeted offense, guiding further dataset refinement.
    \item \textbf{Confidence Histogram (\texttt{confidence\_histogram\_unseen.png}):} This chart mapped the statistical probability distributions of the AI's predictions. It demonstrated that the model was highly confident (scoring above 95\%) on clear hate speech, allowing the frontend to apply extreme blurring with certainty.
\end{enumerate}

\begin{figure}[htbp]
    \centering
    \safeincludegraphics[width=0.8\textwidth]{prism-uploads/confusion_matrix_2.png}
    \caption{Confusion Matrix displaying classification accuracy across all buckets.}
    \label{fig:confusion}
\end{figure}

\begin{figure}[htbp]
    \centering
    \safeincludegraphics[width=0.8\textwidth]{prism-uploads/confidence_histogram_2.png}
    \caption{Confidence Histogram showing probability distributions of predictions.}
    \label{fig:confidence}
\end{figure}

\chapter{CONCLUSION AND FUTURE SCOPE}

\section{Conclusion}
The Safe Chat System represents a massive leap forward in the automated moderation of regional, code-mixed digital communication. By meticulously engineering a robust dataset and combining a mathematical Regex/Cosine-Similarity firewall with an active Transliteration Engine, the system successfully bridges the gap between chaotic human typing and structured machine learning. 

The fine-tuned MuRIL AI model demonstrates exceptional capability in understanding the deep semantic context of Manglish. Coupled with the live continuous-feedback loop, the Instant Whitelist Cache, and the automated retraining pipeline, the platform is not merely a static filter, but a highly adaptive, self-learning ecosystem. The inclusion of the Web Speech API and the asynchronous polling architecture ensures that the platform delivers a modern, highly responsive, and accessible user experience.

\section{Future Scope}
While the current system is highly capable, several avenues for future enhancement exist:
\begin{itemize}
    \item \textbf{WebSocket Integration:} Upgrading the Live Polling architecture from an HTTP \texttt{fetch} loop to a true WebSocket-based stream (e.g., utilizing \texttt{Flask-SocketIO}) would eliminate polling overhead and provide true zero-latency, bi-directional communication.
    \item \textbf{Cloud Deployment and API Scalability:} Deploying the fine-tuned model to a dedicated cloud cluster (e.g., AWS SageMaker or Google Cloud AI) would allow the moderation engine to be offered as a standalone REST API.
\end{itemize}

\clearpage

\appendix
\chapter{CORE IMPLEMENTATION EXCERPT}

This appendix includes only the essential backend code required to understand the Safe Chat runtime workflow. Non-essential setup code, local environment configuration, auxiliary dataset scripts, metric utilities, downloader scripts, and full source listings are omitted to keep the appendix concise and focused.

The excerpt below captures the main implementation path from \texttt{prism-uploads/chat\_server.py}: text normalization, transformer-based classification, live message storage, and feedback persistence for active learning.

\section{Core Backend Moderation Flow}
\begin{lstlisting}[language=Python]
"""
Abridged Safe Chat backend excerpt.
Shows only the core runtime pipeline: normalize text, classify it,
store chat messages, and persist administrator feedback.
"""

import csv
import os
import re
import uuid
from datetime import datetime

from flask import Flask, jsonify, request
from indic_transliteration import sanscript
from sklearn.feature_extraction.text import TfidfVectorizer
from sklearn.metrics.pairwise import cosine_similarity
from transformers import pipeline

app = Flask(__name__)

DATASET_PATH = os.path.join(os.path.dirname(__file__), "mal_full_offensive_train.csv")
FEEDBACK_LOG = os.path.join(os.path.dirname(__file__), "feedback_log.csv")
MODEL_DIR = os.path.join(os.path.dirname(__file__), "finetuned_model")
USER_CHAT_DATASET = os.path.join(os.path.dirname(__file__), "user_chat_dataset.csv")

classifier = pipeline("text-classification", model=MODEL_DIR, top_k=None)
messages = {}
message_order = []
USER_OVERRIDES = {}

SAFE_LABELS = {"Not_offensive", "Not_in_intended_language"}
LABEL_META = {
    "Not_offensive": {"display": "Safe", "severity": "safe"},
    "Not_in_intended_language": {"display": "Not Malayalam", "severity": "safe"},
    "Off_target_group": {"display": "Group Hate Speech", "severity": "high"},
    "Profanity": {"display": "Profanity / Vulgarity", "severity": "high"},
    "Off_target_ind": {"display": "Personal Attack", "severity": "high"},
}

BASE_SLURS = ["..."]
SAFE_EXCEPTIONS = {"kashuvandi", "cashew", "apple", "pooryum"}
_vectorizer = TfidfVectorizer(analyzer="char_wb", ngram_range=(2, 3))
_vectorizer.fit(BASE_SLURS)
_base_vectors = _vectorizer.transform(BASE_SLURS)


def apply_cosine_similarity(text, threshold=0.85):
    processed_words = []
    for word in text.split():
        clean_word = re.sub(r"(.)\1{2,}", r"\1\1", word.lower())
        if clean_word in SAFE_EXCEPTIONS:
            processed_words.append(word)
            continue

        word_vector = _vectorizer.transform([clean_word])
        scores = cosine_similarity(word_vector, _base_vectors)[0]
        best_index = scores.argmax()
        processed_words.append(BASE_SLURS[best_index] if scores[best_index] >= threshold else word)

    return " ".join(processed_words)


def preprocess_text(text):
    text = sanscript.transliterate(text, sanscript.MALAYALAM, sanscript.OPTITRANS)
    text = text.lower()

    replacements = {
        r"\bchettan\b": "brother",
        r"\bchetta\b": "brother",
        r"\bpoda\b": "friend",
        r"\bpodi\b": "friend",
        r"\bpoo+\b": "flower",
        r"\bpoor(?:i|yum)\s+(?:and\s+)?(?:curry|bhaji|masala|kazhikan)(?:um)?\b": "food",
        r"\bninte\b": "your",
        r"\bninak(?:k)?\b": "you",
        r"\bda\b": "friend",
    }
    for pattern, replacement in replacements.items():
        text = re.sub(pattern, replacement, text)

    return apply_cosine_similarity(text)


def classify(text, is_audio=False):
    raw_text = text.strip().lower()
    if raw_text in USER_OVERRIDES:
        label = USER_OVERRIDES[raw_text]
        is_offensive = label != "Not_offensive"
        meta = LABEL_META.get(label, {"display": label, "severity": "high"})
        return {
            "label": label,
            "label_display": meta["display"],
            "offensive_score": 1.0 if is_offensive else 0.0,
            "safe_score": 0.0 if is_offensive else 1.0,
            "is_offensive": is_offensive,
            "bucket": meta["severity"],
            "all_scores": {label: 1.0},
        }

    processed_text = text.strip() if is_audio else preprocess_text(text)
    results = classifier(processed_text)[0]
    safe_score = sum(item["score"] for item in results if item["label"] in SAFE_LABELS)
    offensive_score = 1.0 - safe_score
    top = max(results, key=lambda item: item["score"])

    if offensive_score <= 0.80:
        bucket = "safe"
    elif offensive_score < 0.95:
        bucket = "high"
    else:
        bucket = "extreme"

    meta = LABEL_META.get(top["label"], {"display": top["label"], "severity": bucket})
    return {
        "label": top["label"],
        "label_display": meta["display"],
        "offensive_score": round(offensive_score, 4),
        "safe_score": round(safe_score, 4),
        "is_offensive": offensive_score > 0.80,
        "bucket": bucket,
        "all_scores": {item["label"]: round(item["score"], 4) for item in results},
    }


@app.route("/api/send", methods=["POST"])
def send_message():
    data = request.get_json(force=True)
    text = data.get("text", "").strip()
    if not text:
        return jsonify({"error": "empty text"}), 400

    result = classify(text, is_audio=data.get("is_audio", False))
    message = {
        "id": str(uuid.uuid4()),
        "text": text,
        "sender": data.get("sender", "me"),
        "username": data.get("username", "User"),
        "timestamp": datetime.now().strftime("%H:%M"),
        "date": datetime.now().isoformat(),
        "revealed": False,
        "reported": False,
        **result,
    }
    messages[message["id"]] = message
    message_order.append(message["id"])

    with open(USER_CHAT_DATASET, "a", newline="", encoding="utf-8") as file:
        csv.writer(file, delimiter="\t").writerow([text, result["label"], "user_chat"])

    return jsonify(message), 200


@app.route("/api/messages", methods=["GET"])
def get_messages():
    return jsonify([messages[message_id] for message_id in message_order]), 200


@app.route("/api/feedback", methods=["POST"])
def feedback():
    data = request.get_json(force=True)
    message_id = data.get("message_id")
    correct_label = data.get("correct_label", "Profanity")
    if message_id not in messages:
        return jsonify({"error": "message not found"}), 404

    message = messages[message_id]
    message["reported"] = True
    USER_OVERRIDES[message["text"].strip().lower()] = correct_label

    with open(DATASET_PATH, "a", newline="", encoding="utf-8") as file:
        csv.writer(file, delimiter="\t").writerow([message["text"], correct_label, "feedback"])

    with open(FEEDBACK_LOG, "a", newline="", encoding="utf-8") as file:
        csv.writer(file).writerow([
            datetime.now().isoformat(),
            message_id,
            message["text"],
            correct_label,
            message.get("offensive_score", ""),
            message.get("label", ""),
        ])

    return jsonify({"status": "ok", "message": "feedback saved"}), 200
\end{lstlisting}

\end{document}
